\documentclass[11pt,a4paper]{article}

\usepackage[utf8]{inputenc}
\usepackage[T1]{fontenc}
\usepackage[portuguese]{babel}
\usepackage{amsmath}
\usepackage{amsfonts}
\usepackage{graphicx}
\usepackage{geometry}
\usepackage{xcolor}
\usepackage{booktabs}
\usepackage{hyperref}
\usepackage{listings}

\geometry{left=2.5cm, right=2.5cm, top=3cm, bottom=3cm}

% --- Configuração para exibição de código Python ---
\definecolor{codegreen}{rgb}{0,0.6,0}
\definecolor{codegray}{rgb}{0.5,0.5,0.5}
\definecolor{codepurple}{rgb}{0.58,0,0.82}
\definecolor{backcolour}{rgb}{0.95,0.95,0.92}

\lstdefinestyle{mystyle}{
    backgroundcolor=\color{backcolour},   
    commentstyle=\color{codegreen}\bfseries,
    keywordstyle=\color{blue},
    numberstyle=\tiny\color{codegray},
    stringstyle=\color{codepurple},
    basicstyle=\ttfamily\footnotesize,
    breakatwhitespace=false,         
    breaklines=true,                 
    captionpos=b,                    
    keepspaces=true,                 
    numbers=left,                    
    numbersep=5pt,                  
    showspaces=false,                
    showstringspaces=false,
    showtabs=false,                  
    tabsize=4
}
\lstset{style=mystyle}

% --- Identificação da Equipe ---
\title{\textbf{Relatório Técnico - Crédito 2}\\ \Large{Implementação do FCM e Obtenção Empírica da Função Tempo}}
\author{
    \textbf{Equipe XX} - [Nome da sua Restrição Arquitetural] \\
    Membro 1: [Nome Completo e Matrícula] \\
    Membro 2: [Nome Completo e Matrícula] \\
    \textit{Universidade do Estado da Bahia (UNEB) - DCET I}
}
\date{\today}

\begin{document}

\maketitle

\section{Introdução e Escopo}
Este documento apresenta os resultados do Crédito 2 da disciplina de Engenharia de Programas. O objetivo é descrever a implementação desenvolvida \textit{from scratch} do algoritmo \textbf{Fuzzy C-Means (FCM)}, respeitando rigorosamente a restrição arquitetural sorteada pela equipe, e apresentar o levantamento do Tempo Empírico ($T_{exp}$) em função do número de pixels ($n$) e clusters ($K$).

\section{Análise da Restrição Arquitetural}
\textbf{Restrição Sorteada:} [INSERIR NOME DA RESTRIÇÃO, ex: Missão "Fuga da FPU" (Fixed-Point Math)]

% INSTRUÇÃO AOS ALUNOS:
% Expliquem detalhadamente COMO vocês lidaram com a restrição arquitetural no código. 
% Quais foram os desafios técnicos de hardware impostos por ela?
% Como converteram a lógica nativa do FCM para obedecer à regra? 
% Exemplo: Se for Fixed-Point, expliquem a escala usada e os bitwise shifts aplicados. 
% Se for Loop Unrolling, expliquem como desenrolaram e lidaram com o "resto" da divisão do laço.
[Escreva aqui a fundamentação teórica de como a restrição foi superada do ponto de vista de Engenharia de Software e Microarquitetura...]

\section{Implementação do Núcleo FCM}
O código fonte completo, contendo os scripts de validação e a implementação integral, encontra-se versionado no seguinte repositório público: \\
\textbf{Link do Repositório:} \href{https://github.com/sua-equipe/projeto-foca}{[INSERIR URL DO GITHUB AQUI]}

Abaixo, destaca-se o \textit{snippet} do laço principal do algoritmo, onde a restrição arquitetural atua sobre as primitivas de processamento e acesso à memória na fase de associação (cálculo das distâncias e atualização da matriz de pertinência).

% INSTRUÇÃO AOS ALUNOS:
% Colem abaixo APENAS o coração computacional do algoritmo (os laços for iterativos e o cálculo interno de pertinência/distância).
\begin{lstlisting}[language=Python, caption=Núcleo computacional do FCM adaptado à restrição.]
# [INSERIR O CÓDIGO FONTE AQUI]
# Exemplo genérico estrutural:
for i in range(N):
    for j in range(K):
        # Implementação autoral da equipe respeitando a restrição
        pass 
\end{lstlisting}

\section{Metodologia Experimental e Estabilidade}
As medições empíricas foram executadas utilizando o script \texttt{experimentador.py}. Para atestar a confiabilidade do hardware e minimizar os ruídos do Sistema Operacional, foram realizadas $R = 30$ repetições para cada combinação de $(n, K)$. 

A Tabela \ref{tab:estabilidade} comprova o rigor estatístico da bancada, garantindo que o Coeficiente de Variação (CV) permaneceu dentro da zona de aceitação ($CV < 0.15$).

% INSTRUÇÃO AOS ALUNOS:
% Preencham a tabela abaixo resumindo as saídas do Pandas DataFrame gerado pelo seu script.
\begin{table}[h!]
\centering
\caption{Validação de Estabilidade Empírica (Amostra Resumida)}
\label{tab:estabilidade}
\begin{tabular}{@{}ccccc@{}}
\toprule
\textbf{Tamanho da Imagem ($n$)} & \textbf{Clusters ($K$)} & \textbf{Tempo Médio ($s$)} & \textbf{Desvio Padrão} & \textbf{CV} \\ \midrule
$100 \times 100$ (10000) & 2 & [0.000] & [0.000] & [< 0.15] \\
$100 \times 100$ (10000) & 3 & [0.000] & [0.000] & [< 0.15] \\
... & ... & ... & ... & ... \\
$500 \times 500$ (250000) & 4 & [0.000] & [0.000] & [< 0.15] \\ \bottomrule
\end{tabular}
\end{table}

\section{Resultados Empíricos ($T_{exp}$)}

A partir da filtragem de estabilidade, o algoritmo de \textit{curve fitting} calibrou a função empírica do sistema. A equação polinomial identificada pela bancada real de testes foi:

% INSTRUÇÃO AOS ALUNOS:
% Colem aqui a equação gerada no terminal pela função exibir_equacao_ajustada().
\begin{equation}
    T_{exp}(n, K) = (\alpha) \cdot nK + (\beta) \cdot n + (\gamma) \cdot K + (\delta)
\end{equation}

Para validação visual, as Figuras \ref{fig:2d} e \ref{fig:3d} apresentam, respectivamente, os cortes bidimensionais e a superfície de resposta tridimensional projetada sobre as aferições reais (pontos de amostragem).

% INSTRUÇÃO AOS ALUNOS:
% Certifiquem-se de que a imagem 'relatorio_complexidade.png' está na mesma pasta do arquivo .tex principal.
\begin{figure}[h!]
    \centering
    % \includegraphics[width=1\textwidth]{relatorio_complexidade.png}
    \vspace{4cm} % Remover esta linha e descomentar a linha de cima ao inserir a imagem real
    \caption{Tempo empírico segmentado: $T \times n$ (esquerda) e $T \times K$ (direita).}
    \label{fig:2d}
\end{figure}

% INSTRUÇÃO AOS ALUNOS:
% Certifiquem-se de que a imagem 'superficie_3d_foca.png' está na mesma pasta.
\begin{figure}[h!]
    \centering
    % \includegraphics[width=0.8\textwidth]{superficie_3d_foca.png}
    \vspace{5cm} % Remover esta linha e descomentar a linha de cima ao inserir a imagem real
    \caption{Superfície de Tempo 3D sobreposta aos pontos de calibração empírica.}
    \label{fig:3d}
\end{figure}

\section{Conclusão Parcial}
Os dados atestam que a implementação adaptada do algoritmo FCM possui complexidade de ordem $\mathcal{O}(nK)$, com os gargalos microarquiteturais refletidos primariamente no coeficiente dominante $\alpha$ da medição. Este modelo empírico servirá como base absoluta (\textit{ground truth}) para o confronto teórico exigido no Crédito 3 utilizando a contagem F.O.C.A.

\end{document}